% ============================================================
%  F. C. Sibbern, Bemærkninger og Undersøgelser,
%  fornemmelig betreffende Hegels Philosophie,
%  betragtet i Forhold til vor Tid
%  Kjøbenhavn: C. A. Reitzel (Bianco Lunos Bogtrykkeri), 1838.
%  TRANSCRIPTION (Danish)
%  Source: KB scan (hegels-philosophie.pdf). PDF page = printed page + 11.
% ============================================================
\documentclass[12pt,a4paper]{book}
\usepackage[utf8]{inputenc}
\usepackage[T1]{fontenc}
\usepackage{libertinus}
\usepackage{libertinust1math}
\usepackage{textalpha}   % Greek words typed directly
\usepackage{graphicx}    % for \reflectbox (the old Danish »ɔ:« id-est mark)
\DeclareUnicodeCharacter{0254}{\reflectbox{c}}  % ɔ (U+0254) = reversed c
\usepackage[danish]{babel}
\usepackage[protrusion=true,expansion=true]{microtype}
\usepackage{setspace}
\usepackage[margin=1.2in]{geometry}
\usepackage{enumitem}
\usepackage{fancyhdr}
\setlength{\headheight}{28pt}
\addtolength{\topmargin}{-13.5pt}
\usepackage{hyperref}

\hypersetup{
  pdftitle={Bemærkninger og Undersøgelser, fornemmelig betreffende Hegels Philosophie},
  pdfauthor={F. C. Sibbern},
  hidelinks
}

% Running heads reproduce the printed "Overskrivter" (set per section via \markboth).
\pagestyle{fancy}
\fancyhf{}
\fancyhead[LE,RO]{\thepage}
\fancyhead[CE]{\textit{\leftmark}}
\fancyhead[CO]{\textit{\rightmark}}

\setstretch{1.3}

\title{\textbf{Bemærkninger og Undersøgelser,}\\[6pt]
  \large fornemmelig betreffende\\[6pt]
  \Large\textbf{Hegels Philosophie,}\\[6pt]
  \large betragtet i Forhold til vor Tid,\\[10pt]
  \normalsize særskilt aftrykte af en i Maanedsskrift for Litteratur
  10de Aargang (1838) indrykket Recension over Prof.\ Heibergs Perseus No.\ 1}
\author{Frederik Christian Sibbern}
\date{Hos Universitetsboghandler C.\ A.\ Reitzel.\\
  Trykt i Bianco Lunos Bogtrykkeri.\\[4pt]
  1838.}

\begin{document}
\maketitle
\thispagestyle{empty}
\newpage

%% ============================================================
%%  FORTALE (printed pp. III--IV)
%% ============================================================
\chapter*{Fortale}
\markboth{Fortale}{Fortale}
\thispagestyle{plain}

Allerede i nogen Tid havde jeg næret den Tanke, engang at yttre mig
offentlig over den Hegelske Philosophie, da Prof.\ Heibergs Perseus
udkom, og ved sin første Artikel: Recensionen over Dr.\ Rothes Skrivt
om Treenighed og Forsoning, gav mig Anledning og Ansporelse til at
udføre hiin Tanke. Saaledes opstod den Recension over Prof.\ Heibergs
Perseus, som begyndtes i Aprilheftet af „Maanedsskrift for Litteratur“,
og er bleven fortsat igjennem de følgende Hefter, og endnu i det
Øieblik, jeg skriver Dette, ikke er heelt aftrykt. For de Læseres Skyld,
som ei holde Maanedsskrivtet, især for dem deriblandt, der høre til mine
akademiske Tilhørere, bad jeg Forlæggeren, Hr.\ Universitetsboghandler
Reitzel, foranstalte særskilte Aftryk af den Deel af Recensionen, som
directe har Hegels Philosophie til Gjenstand. Imidlertid er dog ogsaa
Begyndelsen af min Recension, der allene har med Prof.\ Heiberg og hans
Perseus at gjøre, her bleven aftrykket, fordi jeg forestillede mig, at
den ikke uden stor Ubeqvemmelighed kunde tages bort. Den sidste Halvdeel
af Recensionen derimod, der ogsaa directe blot angaaer Prof.\ Heiberg og
den af ham hverken billigt eller retfærdigt eller endog kun nogenlunde
kritiskt behandlede Rotheske Bog, er ikke bleven optagen i disse
særskilte Aftryk. Men ogsaa den indeholder dog en Deel Bidrag til
Opklaringen af enkelte vigtige Puncter af Hegels Lære, (f.\ Ex.\ Hegels
Begreb om det af ham saakaldede slette Uendelige og hans Lære angaaende
Guds Ubegribelighed), hvilke jeg nok vilde ønsket at have kunnet give
Plads i nærværende særskilte Aftryk, dersom det kunde have ladet sig
gjøre, uden at tage en Deel med, som slet ikke hørte herhen, samt at
gjøre alt for store Fordringer til Forlæggeren. Dem maa jeg da nu bede
Læseren at ville søge i selve Maanedsskrivtet, hvilket da ei heller
Forlæggeren kan ønske gjort overflødigt for Nogen ved at tilstede
særskilte Aftryk i et for stort Omfang; jeg har selv ei turdet
tilsidesætte hans og selve Maanedsskrivtets Interesse i saa Henseende.

Ovenstaaende havde jeg ved disse, ikke i noget stort Antal tagne,
særskilte Aftryk forud at bemærke. Endnu maa jeg gjøre opmærksom paa
Rettelserne, af hvilke de fleste kun vedkomme disse særskilte Aftryk,
idet flere af dem betreffe Henviisninger til Paginer, der paa de første
Ark af Aftrykkene svare til Maanedsskrivtets Paginer, men ei til
Aftrykkenes, andre angaae Overskrivterne over Siderne, hvilke formedelst
Columnernes Omflytning ved de særskilte Aftryk nogle Gange ere komne til
at staae urigtig. En Oversigt over Indholdet føies til, saaledes som den
og i Maanedsskrivtet fandtes paa Omslaget af hvert Hefte og er optagen i
sammes Register.

\medskip
\noindent Kjøbenhavn 31te August 1838.\hfill \textbf{Sibbern.}

%% ------------------------------------------------------------
%%  RETTELSER (errata, printed p. IV)
%% ------------------------------------------------------------
\bigskip
\begin{center}
  \textbf{\large Rettelser.}
\end{center}

\noindent\emph{I Citationerne:}
\begin{itemize}[leftmargin=3.2em,itemsep=0pt]
  \item[Pag.\ 20] nederste Linie: 294-97 læs: 12-15.
  \item[--- 41] L.\ 4 istedetf.\ 301 sæt: 19.
  \item[--- 43] L.\ 7 fra neden: istedetf.\ 298 sæt: 16.
  \item[--- 62] L.\ 21 istedetf.\ 328 sæt: 46.
  \item[--- 95] L.\ 6 --- 296 --- 14.
  \item[--- 111] nederste L.\ istedetf.\ 311 sæt: 29.
\end{itemize}

\noindent\emph{I Overskrivterne:}
\begin{itemize}[leftmargin=3.2em,itemsep=0pt]
  \item[--- 91] istedetfor Overskrivten: Om den formale Logik, sæt:
    \textit{principium indiscernibilium}, hvilke to Ord derimod Pag.\ 93
    i Overskrivten udslettes.
  \item[--- 95] skal Overskrivten forandres til: \emph{Tænkningens Suprematie.}
  \item[--- 99] skal Overskrivten forandres til: \emph{Tre Slags Erkjenden;}
  \item[--- 111] ligeledes.
\end{itemize}

\noindent\emph{Desuden er Følgende at rette:}
\begin{itemize}[leftmargin=3.2em,itemsep=0pt]
  \item[Pag.\ 2] L.\ 6 fra neden: Skikkelsen læs: \emph{Skikkelse}; bundne l.\ \emph{lænkede.}
  \item[--- 21] L.\ 6 og 5 fra neden: Mere, læs: \emph{mere.}
  \item[--- 32] L.\ 6: Sufficence, l.\ \emph{Suffisance.}
  \item[--- 39] L.\ 13: kan nu, læs: \emph{kan nu ogsaa.}
  \item[--- 50] L.\ 5 i Noten: Ordet \emph{givet} skal udmærkes.
  \item[--- 52] L.\ 8 sættes Semikolon istedetf.\ Kolon.
  \item[--- 53] L.\ 3 og 4 baade være, l.\ \emph{være baade.}
  \item[--- ---] L.\ 3 fra neden: Eet, læs: \emph{Et.}
\end{itemize}

\cleardoublepage

%% ============================================================
%%  HOVEDTEKST — begins printed p. 1 (PDF 12)
%%  Running review of Heiberg's Perseus No. 1
%% ============================================================
\pagenumbering{arabic}
\markboth{Heibergs Perseus.}{Heibergs Perseus.}

Var det ikke Professor Heiberg, der traadte frem som Udgiver af den
Journal for den speculative Idee, vi her have for os, saa vilde man
maaskee kunne finde Grund til at bemærke, at der, i Conseqvents med hvad
strax i Begyndelsen af den første Hovedartikel er sagt om Philosophiens
Lod og Skjæbne hos os, ikke burde være bleven brugt en saa megetsigende
Titel. Men Prof.\ Heiberg har kunnet være vis paa, at om et philosophiskt
Tidsskrivt, han udgav, vilde Enhver strax forudsætte, at han deri vilde
lade sig det være Hovedsagen at popularisere den speculative Erkjendelse
(see Pag.\ \textit{XI} i Fortalen), og at gjøre den interessant, uden at
føre sine Læsere dybere ind i den, eller fordre mere Anstrengelse af dem,
end den Danske, med sit deels lette, deels magelige Sind, nok kan
overkomme, og at Forfriskninger af munter Natur ei vilde mangle. Han selv
endog har --- med alt det at han minder om, at han ei bør nedværdige sit
Formaal ved „at gaae uden om Speculationen“, eller for Popularitetens
Skyld „opoffre eller endog blot tilhylle“ hvad „der er Livsprincipet i
Videnskaben“ --- stemmet de Fordringer, han vil tilfredsstille, temmelig
meget ned, ved at betegne de Læsere, paa hvis Understøttelse hans
Foretagendes Fremgang vil komme til at beroe, som „Læsere, der maatte ane
Nødvendigheden af, at den humane Dannelse, de have erhvervet sig, nu ogsaa
engang udbreder sig til deres Tænkning“
(Pag.\ \textit{XIII}). Jeg veed meget vel, at Mange ville forstaae disse
Ord, især ved at sammenholde dem med hvad Udgiveren længere hen
(Pag.\ 8) siger om Philosophiens Tilstand hos os, som indeholdende et
Udbrud af en Forestilling, der er grundet i en temmelig fornem Nedseen
paa Læserne, hvilke her stilles saa lavt, at de først skulle til at ane,
at de dog engang maae til at lade deres humane Dannelse udbrede sig til
deres Tænkning. Thi her staaer ikke engang: til deres Tænknings
philosophiske Uddannelse; nei her staaer simpelthen: „til deres
Tænkning.“ Men man skal ei kunne nægte mig, at disse Ord ogsaa meget vel
kunne udlægges som Udtryk af Beskedenhed, idet Udgiveren, som forresten
vist ikke vil have Noget imod, at finde Læsere, hvis humane Dannelse for
længe siden er begyndt at udbrede sig til deres philosophiske Tænkning,
dog, kjendende sin Formaaens Natur og Omkreds, ei vil træde frem med den
udtrykkeligen udtalte Erklæring, at han vil skrive for saadanne.
Forfatteren vil virke vækkende, kan jeg troe; men de Vakte behøve ei at
vækkes, og at tilfredsstille disse vil han mueligen ei paatage sig, om
end vel nok ogsaa af saadanne vente sin Journal læst med Interesse.

At han kalder sit Tidsskrivt Perseus, har sin Grund i den Udtydning, han
giver Mythen om denne Helt. „Af Gorgonens (Medusas) Blod, den ophævede
Realismes forsvindende Livskilde“ --- saa yttrer han sig i den Forklaring,
han giver --- „fremsprang Pegasus, den poetiske Ganger, paa hvis bevingede
Ryg Helten befriede den hidtil bundne, for Naturmagten priisgivne Idee,
under Skikkelsen af den til Klippen bundne Andromeda.“ --- „Vel betyder“
--- saa hedder det videre --- „Perseus paa Dansk Ødelæggeren; men hans
Ødelæggen var ikke Tidens og Dødens absolut negative, men den heroiske,
hvoraf en ny og ædlere Tilstand fremgaaer.“ At her ved Perseus især maa
tænkes paa Dialektiken, er vel klart. Lignelsen er heldig udfunden og vel
gjennemført. Det forstaaer sig, at Ordet Realisme her maa tages i sin
vulgære Betydning.

Saa Meget angaaende Fortalen, paa Eet nær, som jeg ei kan lade være at
omtale. Det vedrører Noget, hvorom jeg efter Göthe kan sige: „Das höre
ich hundertmal wiederholen, Ich fluche darauf, aber
verstohlen.“*)\footnote{See Göthes Schriften, Th.\ 3, under den Titel:
Allerdings.} Forunderligt, at endog Prof.\ Heiberg, som dog ellers med
Hegel ynder Treheden, og overalt ellers snarere viser sig tilbøielig til
at gribe efter Nyt, end at ynde gammel Suurdeig, i eet Punkt bliver ved en
Tohed, der dog nu vel maatte have udtjent. Jeg mener den æsthetiske
Distinction imellem Form og Indhold. Aabenbar ere der i den Henseende, i
hvilken denne Distinction er gjort, tre Ting at adskille i ethvert
Konstværk (ligesom i enhver levende Organismus). Det Første er det
Stofartige, det sig til Bearbeidelse frembydende Givne; dertil træder det
bearbeidende og sig i dette Givne indarbeidende livfulde Princip, Ideen,
det egentlige Poetiske, svarende til Det, som i philosophisk Henseende er
det Speculative, med sin hele Gehalt; og endeligen have vi for det Tredie
at lægge Mærke til dettes Indorganisation i hiint, hvilken gaaer i Eet med
dets Udfoldning eller Udorganisation, som er Det, man ved Udtrykket: Form
betegner. Taler man om Behandling, saa maa man dog tilstaae, at her maa
være at skjelne imellem det Behandlede, Behandlingen, og den sig i og
igjennem Behandlingen udtalende Aand. Nævner man Form, saa maa det vel
være klart, at her maa være Noget, om hvis Formelse der er Tale, men at
begge forudsætte Noget, som gjennem Formen træder ud i det
Formede.**)\footnote{Blandt de mange ypperlige „Maximen und Reflexionen“
af Göthe, som læses i det 9de Bind af hans efterladte Værker, findes
følgende, hvori den her opstillede Trehed er udtalt: „Den Stoff sieht
jedermann vor sich, den Gehalt findet nur der, der etwas dazu zu thun hat,
und die Form ist ein Geheimniß den Meisten.“} Om man
just hver Gang kan sondre disse Tre saaledes ad, at man kan eftervise Det
og Det som det blot Stofartige, Det og Det som henhørende blot til
Formationen eller Organisationen, og Det og Det som Udtrykket af den Aand,
der gjennem denne Formation træder ud i Værket, derpaa kommer det ikke an.
De Tre kunne paa hvert Punct uadskilleligen gribe i hinanden, ja de maae
det, uden at de derfor høre op at være at adskille for Betragtningen, som
Tre. --- Sujet, Idee, Plan: med disse tre Udtryk ere de Tre ofte nævnede;
de vise os Treheden som den, der allerede i Konstværkets Opstaaen fra den
ideelle Side gjør sig gjeldende. Man kan nu vistnok sige, at det først er
den opstaaende Idee, der forvandler en Opgave til et Sujet i egentlig
Forstand, skjønt dog allerede Opgaverne ofte med Grund kaldes Sujetter,
som de, hvormed Tænkningen beskjæftiger sig; men især kan man sige, at,
hvor Opgaver ei gives udenfra, men hvor først Aandens frie Virken ladet
Noget blive til et Sujet, der er det først den opstaaende Idee, som gjør,
at Sujettet overhovedet bliver til, og i denne Idee ligger da strax
allerede Grundspiren til Planen. Men Dette kan ingenlunde hindre at see
den heri tilstedeværende Trehed. Under Udførelsen, d.\ e.\ under Ideens
Arbeiden i Digtergemytet, kommer det altid an paa, hvad der kommer fra
Aandens Side, fra den speculative Verdensbetragtnings, fra det poetiske
Bliks, det poetiske Gemyts, fra Lunets, Humorens, Ironiens, fra den dybere
Menneskefølelses Side, o.\ s.\ v.; --- dernæst kommer det an paa, om Det,
der især skal komme fra Konstens Side, træder til, om der er en sand
sammenhængende saakaldet Plan, en sand Organisation og Eenhed i den hele
Formation, som baade consoliderer og udfolder; og endeligen, om de begge,
Aand og Plan, vise sig mægtige over det Indbegreb af Enkeltheder, der, saa
at sige, skulle overmandes ved Formen, eller assimileres ved den og vorde
indorganiserede saaledes i den, at Ideen i og gjennem dem kommer til at
træde ud baade i sin Dybde og med Eenhed og Heelhed. Naturligviis
repeterer Dette sig paa ethvert Punct.

Ved nærmere Overveielse opstaaer endog det Spørgsmaal, om her ikke
egentligen turde være en Fiirhed at bemærke, og det netop saa, at jeg
herved meget vel tør bringe Franz Baaders „Vierzahl des Lebens“ i
Erindring. Naar det dog egentligen bliver Idee, Plan og Detail, eller
detailleret Udførelse, vi her have at adskille, som det Grundbevægende med
sin Fylde, dettes centrale Leven i det Værk, der skabes og organiseres, og
dets peripheriske Leven deri i Enkelthedernes hele Spil, saa kan denne
Trias, saaledes tænkt, ei Andet end ligefrem føre til, derfra at adskille
det egentlige blotte Materielle, som det Fjerde; hvorved jeg mener det
forud for denne hele Trias givne Indbegreb af Fremstillingsmidler, hvoraf
den skal berede sig den Heelhed, hvori den skal komme frem, være det
Sproget, Farverne, Tonerne, o.\ s.\ v. Ved det Udtryk Stof, Stofartigt,
tænker Ingen herpaa. Det Stofartige i Drama, Roman, Epos, er den saakaldte
Fabel, som den tørt kan fortælles, eller det saakaldte Æmne, der af
Historien kan være taget, med Alt, hvad der videre af Historien kan
frembyde sig til Optagelse. Musik, Malerie, Lyrik o.\ s.\ v.\ have i noget
Lignende deres Stof. Dette Stofartige er det, der ved og under Formelsen
behandles og beaandes. Adskille vi da i hiin Trias de Tre: Idee, Plan og
Detail, saa maae vi dog, idet vi i Planen see den i Mangfoldigheden sig
etablerende og vedligeholdende Eenhed, men i Detaillet see
Mangfoldighedens Spil, atter skjelne imellem dette Spil og det Materielle,
hvori det spiller, og hvorigjennem det træder ud i sit umiddelbart
anskuelige Udvortes. Man kan og udtale Dette saaledes: Fornuft eller
Ideegehalt træder formedelst Forstand igjennem Imagination og Vid ud i et
udvortes anskueligt Heelt, som dannes ved Hjelp af Sproget, Farverne,
Tonerne o.\ s.\ v. Idee træder planmæssig, organiserende, (paa et Slags i
æsthetisk Forstand systematisk Maade) ud i Mangfoldighedens Spil, men i en
vis ved et givet udvortes Fremstillingsmiddel bestemt Skikkelse. Naar en
Musiker phantaserer genialt paa et Instrument, saa have vi et
Grundbevægende, som det indre Eet og Alt, der fylder ham og fylder os.
Dette træder ud i et Spil af Ideeassociationer --- saaledes kalde vi det
Løb af Forestilling, eller det lille Bevidsthedens Livsløb, der tillige er
Ideens Livsløb, som udgjør den Phantasering, vi høre. Men heri er heelt
igjennem et Regulerende, og Classiciteten beroer paa dette Regulerendes
gjennemtrængende, Alt til Eenhed samlende og sammenholdende Magt. Men denne
Gang vare det nu Tonerne, som afgave det Indbegreb af Fremtrædelsesmidler,
hvorigjennem hiin Trehed traadte ud. Der kunde være phantaseret paa
lignende Maade gjennem Sproget af en genial Improvisator, eller af en
Taler, som udgjød sin Tankefylde i ubunden Stil. Den samme Trias havde da
virket og skabt gjennem et et andet Materiale.% [sic: printed "et et andet" — line-break dittography in the 1838 print]
Man seer iøvrigt dette blotte Materiales væsentlige Betydenhed, eftersom
dog saadan Skabens hele Art og Maade i Meget bestemmes ved dette
Materiales Beskaffenhed, og det endog saa, at, naar det er Sproget, vil
det ingenlunde blive uden Indflydelse paa hiin indre ideelle Trias,
hvilket Sprog det er, især naar det ere Sprog af stor Forskjellighed, om
hvilke der er Tale. Landskabsmaleren har sit Stof i æsthetisk Betydning i
Luft, Vand, Skov o.\ s.\ v., og ikke i Palettens Farver. Men nu har han
dog ogsaa disse sine Farver at tage Hensyn til, og ogsaa de ville studeres.

\begin{center}---\end{center}

Vi gaae nu fra Fortalen til Journalen selv. Den indeholder tre
Hovedstykker: 1) en Recension af Udgiveren over Dr.\ Rothes, i Anledning
af Reformationsfesten i Efteraaret 1836, som et „speculativt Forsøg“
udgivne Skrivt om Treenighed og Forsoning; 2) en Artikel om Ideen af
Faust, af Lic.\ theol.\ Martensen, med Hensyn til en af en tydsk Digter
(Lenau) leveret Digtning, som har Faust til Gjenstand; 3) en Recension af
Prof.\ Heiberg over Hertz's Svend Dyrings Huus, tilligemed en æsthetisk
Betragtning af de danske Kæmpeviser. --- Man har altsaa her modtaget tre
Recensioner, skjønt alle tre sigtende til at præstere Mere, end man af
Recensioner kan vente, kort Artikler i Lighed med dem, Maanedsskrivtet nu
og da leverer. Man kan da ei recensere nærværende Bog, uden at recensere
Recensioner. Men at blandt andet Meget, som vorder recenseret, ogsaa
engang Recensionerne selv komme for Touren til at recenseres, er vel ganske
i sin Orden. Indirecte tør det vel og mangen Gang være skeet. At det saa
sjeldent, ja næsten aldrig, er skeet directe --- thi Antikritiker ere endnu
ei Recensioner af Recensioner --- har ikke været til Gavn hverken for
Recensionsvæsenet, eller for de Forfattere, som vorde recenserede. Der er
en forunderlig Ulighed deri, at visse Forfattere, nemlig Forfatterne af
Recensioner, endog i saa store Landes Literaturer, som Tydsklands, reent
have været undtagne fra at være Gjenstande for det Recensionsvæsen, som
ellers har søgt at omfatte alle Literaturens Gebeter. Hvad den første af de
her i den for os liggende Bog indeholdte Recensioner angaaer, da har dens
Forfatter ogsaa belavet sig paa, at den vilde kunne blive „Gjenstand for
en ny Recension“ (see Pag.\ 10 øverst).

Til denne Journalens første Artikel gaae vi da nu først, og komme dermed
ind i den hele lille Samling af philosophiske Bemærkninger og
Betragtninger, der hos Prof.\ Heiberg have knyttet sig til adskillige
Steder i Dr.\ Rothes ovennævnte Bog. Dog har vor Forf.\ indskrænket sig
til den Deel af Dr.\ Rothes Afhandling, som angaaer Treenighedslæren; thi
Forsoningslæren berører Prof.\ Heiberg ikke, hvilket allerede den
fortløbende Overskrivt, Side for Side, over Artikelen viser.

Prof.\ Heiberg har her leveret os adskilligt Godt, Adskilligt især, som
vel tør være egnet til, ved det Tiltrækkende tillige, som hans Stiil har,
at lokke Mange hen til de speculative Regioner, og at virke til at
„popularisere den speculative Erkjendelse“. Vor Forf.\ viser sig i flere
af de Brudstykker, som her ere radede sammen, aabenbar paa Ny som en
Skribent, hvem man, til Bedste for den philosophiske Idee-Udviklings
Fremgang her i vort, vistnok ikke meget philosophiske, Fædreneland, maa
ønske Drivt og Understøttelse til at udføre, hvad han, ved at begynde sin
Perseus, har tænkt at ville gjøre. Allerede i Aaret 1825 blev i
Litteraturtidenden yttret om ham, at han turde „være i Stand til at kunne
give fortrinlige Bidrag til at opklare de hegelske Speculationer og at
gjøre dem mere tiltrækkende, end de ere det i Hegels eget knudrede Sprog“,
samt at Hegel turde have faa Tilhængere, „der kunde gjøre det med saa megen
Udvikling af Talent og Originalitet.“ Vel har han nu ikke siden uddannet
sig meget i denne Retning, hvorover vi ogsaa paa Poesiens Vegne have
Aarsag til at være meget fornøiede; ogsaa turde han vel nu ikke mere være
særdeles egnet til at tilfredsstille --- jeg vil ei sige grundige Spørgere
--- jeg vil kun sige: tilfredsstille saadanne tænkende Hoveder, som ved ham
selv kunde vorde vakte til en levende Interesse for sand Philosophie; vel
see vi ham tage sig Sagen temmelig let, idet han paa Rhapsode-Viis bevæger
sig i enkelte af de mere let tilgjængelige og plaisante Partier af den
Hegelske Ideekreds; vel viser han sig, for at jeg med eet Ord skal sige
det, kun som en Dilettant i Philosophien. Men med en saa aandfuld og
livelig Dilettant turde det nok være, at Philosophien hos os en Tid lang
meget vel kan være tjent, idet han meget vel vil kunne være skikket til at
vække tænkende Hoveder til at begjære Mere, end han vil yde. Det Udtryk:
Dilettant, veed jeg, vil støde ham, og naar det ei vil saare ham, saa er
det, fordi det vil opirre ham. Men jeg faaer sige Sagen, som den er. Intet
tænkende Hoved hos os vil kunne undlade at bemærke, hvad jeg har tilladt
mig med rene Ord her at udsige. Kjært skulde det være mig, om jeg herved
kunde incitere vor Forfatter til, næste Gang at vise sig som Den, der skal
ville forsøge, om han ei skal kunne gjøre disse mine Ord til Skamme; ---
det vilde da være en saare ønskelig Virkning af disse mine Ord, for hvis
Skyld allene det var værdt at sætte dem hen. Bedre i alt Fald, langt bedre,
end denne hans Recension er skreven, maa Prof.\ Heiberg være istand til at
skrive, naar han blot faaer Impuls nok til at føle, at han maa det.
Forresten maa jeg, hvad slig Dilettantisme angaaer, bemærke, at en saadan
Dilettant-Existents, som medfører, at man kan være stærkt greben af eet
Punct, uden at see eller ændse, hvad der for Kjenderen ligger lige tæt ved,
men at man derfor ogsaa desto stærkere kan yttre sig som den Grebne, netop
maaskee kan gjøre et liveligt Hoved særdeles skikket til at danne Mediet
for den speculative Idees Popularisering, især med Hensyn til Saadanne, der
skulle til at ane Nødvendigheden af at udstrække deres humane Dannelse
ogsaa til denne Region. Mere end et Slags Dilettantbeskjæftigelse med, og
Dilettantfornøielse ved det Speculative kan da den hele Artikel, vi have
for os, naar undtages hvad der til at opklare Læren om Sjælens Udødelighed
er fremsat, ei være lagt an paa at forskaffe. Men selv som saadan kunde dog
Artikelen være affattet med langt mere Sands for Grundighed, jeg kan sige:
med langt mere Sands for Philosophie, end den er det. Uden at have meget
høie Tanker om Prof.\ Heibergs philosophiske Aandsgaver, kan jeg dog i
Sandhed ikke have saa ringe Tanker om dem, at jeg ikke skulde antage, at
han i denne Artikel ingenlunde har været, hvad han maatte kunne være
(\textit{sibi defuit}).

Forresten har jeg ikke kunnet Andet end med Fornøielse blive vaer, at
Prof.\ Heiberg her viser sig for os som Den, der ikke just er saa fængslet
af det Hegelske System, at han jo, naar de rette Impulser komme, kan baade
bevæge sig nogenlunde frit i de hegelske Anskuelser, saa og begynde at gaae
ud over Hegel. De bedste Partier i hans Recension turde netop være de,
hvori det Hegelske er traadt i Baggrunden, og friere Blik paa Verden vise
sig. Den maatte ogsaa ikke have Øine at see med, eller ikke vide af, hvad
det er paa fri og selvstændig Maade at see sig om i Tilværelsen, og at
blive sig sin egen Tilværs Indhold med nogenlunde Klarhed og Bestemthed
bevidst, som ei kunde blive de store Eensidigheder og de store Huller i den
Hegelske Lære vaer, samt ogsaa dens Utilfredsstillende i Henseende til
Meget af Det, den positivt giver; hvilket da ingenlunde mindst gjelder om
Det, man i psychologisk, moralsk og christelig Henseende finder hos Hegel.
Nu! --- blive alt Dette vaer, det gjør nu rigtignok ikke Prof.\ Heiberg;
men hvor han, glemmende sit Hegelske Systems Holdningspuncter, gaaer frit
efter sin egen Genius's Tilskyndelser, see vi ham dog som Den, der blot
behøver at blive sig selv klar, for at alt hiint Stødende og
Utilfredsstillende i Hegel maa blive ham klart. Men det Hegelske faaer
rigtignok hvert Øieblik det gamle Overherredømme tilbage, og det er egent
nok at see, hvorledes det da strax berøver ham Det, han var kommen til, og
skiller ham ved de Udsigter, der begyndte at aabne sig for ham. Meest
paafaldende har denne Dobbelthed, denne Afhængighed af det Hegelske og
denne friere Gaaen ud derover, været mig, da jeg sammenholdt Det, han
Pag.\ 19--20 fremsætter, hvor han kommer lige hen til den faderlige Gud, og
hvor han kjækt bekjender sig til Det, jeg i at afhandle de Puncter, der
henhøre til Læren om den menneskelige Frihed, oftere har talt om, at
Mennesket lige over for Gud ikke er Intet, men Noget, ja, hvor han endog,
maaskee uden selv at mærke det, fører os lige hen til den Gud, til hvilken
man kan bede --- da jeg, siger jeg, sammenholdt Dette med hvad han selv,
hvor det gjelder at tale logiskt og dialektiskt paa Hegels Viis, siger
angaaende Individualitet og Udødelighed, ja hvad han lige strax foran havde
sagt. Thi her (Pag.\ 18) synes han ikke at ville vide af det sande og
egentlige Objectivitetsforhold til Gud, hvilket dog i Det, hvad han strax
efter, hvor hans egen Aand raader friere (Pag.\ 19--20), udtaler, og som
jeg nu omtalte, aabenbar maa danne Grundlaget, da kun under Forudsætning af
et saadant Objectivitetsforhold, hvori Gud selv har sat sig til
Menneskene, som til sine frie Børn, hiin Faderlighed kan
indtræde.*)\footnote{Det forstaaer sig af sig selv, at jeg paa sit Sted,
naar jeg kommer til at gjennemgaae Prof.\ Heibergs Recension, som
Recension betragtet, vil komme til nærmere at gjennemgaae hvad Prof.\
Heiberg her har udtalt.} Men man maa oftere saaledes med Undren blive vaer,
at Forfatteren, efter at være kommen over Aaen til de friskere og bedre
Kilder, tyer tilbage igjen til sin tilvante Aa, for at blive ved dens
Vande. Virkelig, man seer, at ogsaa for ham vil en anden og rigere
„Verdensanskuelse“ begynde at gaae op, end den, Hegelianismen frembyder.
Jeg har ved disse Ord Hensyn til, hvad Prof.\ Heiberg (Pag.\ 33) siger,
hvor han vil lade Haant ad Professor Poul Møller, fordi denne ei vil blive
staaende ved Hegel, hvorom jeg længere hen skal tale nærmere. Han kan ikke
finde sig i, at Saadanne, som en Tid lang have givet sig ind i Hegels
Philosophie, ville gaae ud over Hegel, fordi de ei mere nøies ved ham, og
han mener, at de ligne Dem, der gaae ud over Aaen for at hente Vand. Hvad
nu Dette angaaer, saa kan vel Intet være naturligere, end at man kan føle
Drivt og Tilskyndelse til at gaae ud over selv den friskt rindende Aa, end
sige den Aa, der forekommer En at begynde at stagnere, for at hente Vand af
hiinsides liggende Kilder, hvorfra maaskee selve Aaen turde have faaet
Noget af sit bedste Vand. Om nu hiin fornævnte den Hegelske Philosophies Aa
indeholder de Livsvande, uden hvilke christelig Philosophie ei vel kan
bygge eller boe, det er jo just Sagen, hvorom det gjelder.

\markboth{Hegel i Forhold til vor Tid.}{Hegel i Forhold til vor Tid.}%
Vi Andre have ei kunnet finde dem, men have i mange Henseender maattet
savne dem. Den Maade, hvorpaa man i den Hegelske Skole har søgt efter
Udødelighedsbegrebet, kan ei være Andet, end et mærkeligt og megetsigende
Tegn i denne Henseende, og jeg anfører det just her i Særdeleshed, fordi
det er i Anledning af nogle Yttringer i Poul Møllers bekjendte, virksomme,
meget læste, og derfor vist mine Læsere i friskt Minde værende Artikel
angaaende Udødeligheden her i Maanedskrivtet, at hiin Tale om at gaae over
Aaen, for at hente Vand, af Prof.\ Heiberg er bleven reist.

Man seer let, at det med Udødeligheden i et heelt, gjennemført philosophiskt
System, til hvis Verdensanskuelse den hører væsentligen med, ikke kan
forholde sig saa, at den mueligen kunde være bleven forglemt af dets
Ophavsmand, og da kunde af Tilhængere føies til og knyttes an eet eller
andet Sted, som Noget, hiin ei var kommen til at omtale. Den udgjør ---
især naar en saadan Philosophie vil sætte sig i et bestemt Forhold til
Christendommen, ja vil indoptage dettes væsentligste Grundlag i sig --- et
ei blot saa væsentligt, men og saa indgribende Punct, at den allerede i et
saadant Systems Anthropologie, end sige i dets Psychologie, og fremdeles i
dets hele Moralphilosophie, maa bestemt komme tilsyne, naar alt Dette skal
ville tilegne sig Navn af speculativt, og naar det endog tager saadanne
Lærdomme med, som den om den dyriske Magnetisme, for at haandhæve, og ved
speculative Anskuelser at oplyse, Gyldigheden af at antage en saadan. Men
især gjelder Dette om en Philosophie, som baade i sin naturphilosophiske og
i sin aandsphilosophiske Deel just især gaaer Tilværelsens opadstigende
Trin og Stadier igjennem, derover endog deels forglemmende, deels
tilsidesættende de allervigtigste Betragtninger med Hensyn til Det, der paa
collateral Maade og i collateral Henseende gjør sig gjeldende i Liv og
Tilvær. Thi, naar i en saadan Philosophie Erkjendelsen af Udødelighedens
Realitet virkelig har Plads, hvorledes kan da den trinmæssige Betragtning
aldeles herske frem deri, uden at Menneskelivet her paa Jorden bliver
fremstillet som et Livets Trin, der viser ud over sig selv til et høiere
Trin paa Individualitetens og Personlighedens, saavelsom paa det i
Personlighederne opgaaende og sig deri indlevende aandelige Riges,
Udviklings-Scala? Hvorledes vil det da kunne undgaae Opmærksomheden, at
Livsrækken i Mennesket først i Udødeligheden finder sin Integration? Dog vi
maae rigtignok bemærke, at, saa meget end den opadstigende Idee-Udviklings\-gang
hos Hegel hersker, er det egentlige Historiske i Liv og Tilværelse dog hos
ham ikke kommen til at blive gjeldende, saaledes som f.\ Ex.\ hos Steffens,
der seer hele Livet historiskt. --- Som det hos et menneskeligt Individ maa
gjøre en Hovedforskjel, hvad enten det tager sit Liv som det, hvis hele
Fylde, Formaal og Sphære er allene at søge paa Jorden (saa at det
forsaavidt om samme maa hedde: \textit{qvod petis hic est}), eller om det
tager Livet som det, der her paa Jorden skal modnes en anden Verden imøde,
saa at det veed sig her paa Jorden endnu kun at befinde sig i Propylæerne,
ei i Templet, saaledes maa det og i et philosophiskt System strax i dets
hele Totalitet vise sig, om Udødeligheden er indoptaget deri. Hos Leibnitz,
Kant, Jacobi, ogsaa Fichte den Ældre, behøve vi ei at spørge. At
Udødelighedstroen er heelt gjennemgribende i Christendommen, er bekjendt
nok. Den ligger ogsaa til Grund for den hele romantiske Poesie. Det hører
til Christendommens store Virkninger, at den har ladet den Tanke blive
mægtig i Bevidsthederne, at vi her endnu kun bære det Uforkrænkelige i
Forkrænkelighed, men at det i os liggende aandelige Foster skal finde den
Dag, da det bryder ud af dette Leie, river sin Navlestreng over, med det
hele derhen hørende Apparat af Organer for Ernæring og Respiration, og
gjennem saadant Brud og saadan Mellembød kommer ud til det Dagens Lys, til
det Element og den friere Verden, som det er bestemt for, hvor det skal
aande og inddrage Næring gjennem andre Organer og i andre Sphærer, end
hidtil. Til de stærkeste, af Aanden meest opfyldte, varmeste og
henrykkelsesfuldeste Steder i det nye Testament høre Pauli begeistrende og
dybt indlysende Ord herom i 1 Cor.\ 15, hvilket jeg ei skulde bemærke,
dersom jeg ei vidste, hvor lidet kjendt det nye Testament, som den hele
Bibel, i vor Tid er hos mange Dannede. At nu alt Herhenhørende i den
Hegelske Lære enten reent fattes, eller træder saa tilbage, at man har
kunnet strides om, om Noget deraf er der, eller ikke, det maa man dog
tilstaae, at være i høieste Grad paafaldende, naar man betænker, at denne
samme Philosophie skal gjelde for at have ført Philosophien lige op til
Christendommens Standpunct, at have gjort Christendommen speculativ, og at
have givet den christelige Tænkning en Grundvold, ved hvilken endog
Theologer anmodes om at holde fast.

Men hvad her angaaende Hegel er bemærket, hænger nøie sammen med det
Utilfredsstillende, der er i Hegels Lære med Hensyn til et andet vigtigt
Punct, der udgjør Grundpunctet i den Objectivitetsidee, som hos Hegel alt
for lidet har faaet Rum. Jeg anseer det for at høre til de første
Hovedpuncter i en Philosophie, der skal bringe Lys ind i Aandens Gebeet, at
det erkjendes, at vi ligesaavel indad til, som udad til, staae i Forhold
til noget Andet, end vi, noget sandt Objectivt, og derved staae i Forhold
til, og Communication med en Region og en Existentsfylde, som har sand
Objectivitet lige over for os, ligesom vi have den lige over for den; saa
at vi altsaa ret egentligen have en indadvendt Sands, hvorved vi sandse det
indenfra Givne og træde i Forhold dertil, hvilket Forhold nu ogsaa har sit
væsentlige Sympathetiske, hvori Toheden som Tohed gjør sig lige saa
gjeldende, som Identiteten, ja hvor denne endnu kun viser sig som den, der
i denne Region først skal komme. Hvilket altsammen vistnok ei lader sig
gjennemføre, uden at man tillægger Endeligheden ogsaa fra Endelighedens
Side, lige over for selve Gud, en Realitet, hvormed mange Yttringer hos
Hegel ingenlunde kunne bestaae.

Er det ikke fremdeles ogsaa i Øine faldende, at, hvis Hegel virkelig var
trængt frem til Christendommens Standpunct i sin Philosophering, og var
kommen nogenlunde ind i dens Region og Fylde, maatte han og i den
menneskelige Følelse og i det menneskelige Hjerte, hvilke han nu behandler
paa en vis forunderlig, foragtelig Maade, have seet noget ganske Andet, end
han nu kan see deri. Overhovedet er Livets Sympathiske ham som ukjendt,
ligesom han ei heller i Aandens Rige retteligen agter paa den Naturgrund,
som ligger overalt, og hvoraf Det, der mødes med det Guddommeliges som
saadant fremtrædende Rige, gaaer frem. Ikke at tale om, at, hvad der i
christelig Forstand kaldes Tro, som levende, hengivelsesfuld, tillidsfuld og
pietetsfuld Griben og Grebethed, slet ikke hos Hegel kommer til retteligen
at gjøre sig gjeldende i den hele Stilling og Betydning, som Troen (hvorved
han kun synes at tænke paa et \textit{credere}, ei en \textit{fides},
πίστις) hos ham har faaet. I hans Yttringer, henhørende til hans
Religionslære, i Slutningen af hans Encyklopædie, og i de hans
Forelæsninger over Religionsphilosophien, der efter hans Død ere udkomne,
findes vel Adskilligt, som, om end kun noget lidet, synes at vidne om, at
saadanne Betragtninger dog ogsaa ere begyndte at gaae op for ham; men hvor
ganske anderledes maatte hans hele Aandsphilosophie, ja selv hans
Naturphilosophie, der nu --- underligt nok --- ei gaaer ud fra den
fuldelige Idee om en \textit{natura naturans}, være kommen til at see ud,
naar Det, som her synes at skimte frem hos ham, var trængt igjennem. Om
flere af disse Puncter, saavelsom om Hegels Trinitetslære, saare lidet
svarende til Christendommens, vil jeg komme til at tale mere længere hen.

Jeg har her berørt adskillige Hovedpuncter hos Hegel, som noksom maae vise,
hvorledes det staaer sig med Hegels Forhold til Christendommen, og tilmed
til de Fordringer, der i en Tid, i hvilken der er kommen nyt friskt Røre i
den christelige Ideekredses Dybder, maae gjøres til den Philosophie, der
skal kunne tilfredsstille Tiden. Men jeg berører her ogsaa Noget, hvori jeg
ei har kunnet Andet, end see talende Beviser for, at Hegel selv, hvis han
nu kunde træde frem i sin Manddoms bedste Kraft, og da havde sin egen hele
Philosophie for sig, umueligen kunde ville blive staaende ved den, men
maatte føle sig dreven til at gaae ud over den paa mange Maader, især om han
tillige kunde træde frem friet for denne fatale og kjedsommelige Polemik, i
hvis Vold vi see ham, og som ei blot har havt en høist forstyrrende
Indvirkning paa hans Systems hele Eurythmie, men endog, ved Siden af hans
haardnakne Fastholden ved visse, hans Dannelsestid tilhørende,
Eensidigheder, synes at have grebet forvirrende ind i hans hele Systems
inderste Constitution. Men nu ei at ville gaae ud over denne Philosophie
--- nu, da den ligger heelt for os, og da der har været Tid nok til
Sundelse og til friere Omblik i Verden --- hvad skal man dertil sige? Paa
mange Maader er det mig klart, at denne hele Lære af en grundig Hegelianer
maatte kunne tages saaledes, at Hegel vilde komme til at staae med
Betydningen af at danne Mellemleddet i et Slags stor Hegelsk Trilogie,
uagtet hans Philosophie i sig selv gaaer ud paa at afslutte en Æra, nemlig
den hele Schellingske Periodes Æra*)\footnote{Jeg undlader ei at bemærke,
hvad der vel ellers maa forstaae sig af sig selv, at jeg, naar jeg her
nævner Schelling, bestandig kun har Schelling, som han da var, for Øie.
Schelling, som han nu er, har jeg her slet ikke at tage Hensyn til, og jeg
har ei heller her kunnet tage Hensyn til ham, som han nu er. Selv har han jo
i mange Aar saa godt som Intet leveret os.}. Thi det maa vel erkjendes, at
den Hegelske Philosophie i Grunden kun gaaer ud paa en Udførelse,
Gjennemførelse og Fuldendelse --- og det endda kun i een særdeles Retning
--- af Det, hvad der i den rige Schellingske Tids Foraarsperiode begyndte at
udfolde sig. I denne ene Retning har nu forresten Hegel gjennemført det med
stor Originalitet, og saaledes, at man vilde lade det Hegelske Genie
vederfares alt for liden Ret, naar man ikke heri vilde erkjende den
grandiose Tænkers mægtige Selvstændighed og store Tankestyrke, samt rige
Fylde af gehaltfuld Bemærkning. Hvortil da kommer, at en Ontologie for denne
Philosophie har Hegel først skabt. Blandt de Lignelser man, deels har
anvendt, deels kan anvende, for at betegne Hegels Forhold til hans nærmeste
Forgjængere (til hvilke Schelling ei ene hører, men den hele Schellingske
Periodes Mænd), turde den være den meest treffende, at han staaer i samme
Forhold til dem, som Spinoza til Cartesius og Cartesianerne, medens der ved
Siden af Spinoza var Plads og Rum nok for en Malebranche, og snart efter
maatte opstaae en Leibnitz, kun at Hegels Værk er saa meget rigere og mere
speculativt tilfredsstillende, som vor Tid er rigere og mere speculativ, end
Spinozas. Men allerede deraf maa, for ei at tale om hine store, Hegels
Person tilhørende Eensidigheder, kunne skjønnes, at Hegel ei kan komme til,
saa stærkt at gribe ind i al Videnskabeligheds Rige, som Schelling gjorde
det, eller som Kant tidligere greb ind, skjønt Hegel har grebet meget mere
ind, end Spinoza gjorde det i sin Tid, eller end Fichte den Ældre i sin Tid.
Philosophien vil især kun i de Perioder kunne gribe bevægende ind i Livet, i
hvilke den undergaaer sine Metamorphoser, og udfolder sig i nye Skikkelser
under en Stræben efter nye Løsninger af de philosophiske Grundproblemer,
idet der paa Ny kommer Røre i disse. Der kom et saadant nyt Røre i dem i den
kantisk-fichtisk-schellingske Tid; men egentligen er intet af dem paa en
gjennemgribende Maade discuteret paa Ny af Hegel, der altfor umiddelbart og
thetiskt har lagt sine Grundanskuelser til Grund, og forresten har ladet
saadanne Problemer, som Spørgsmaalet om Sjælens Forhold til
Legemet*)\footnote{Hvor han berører dette Forhold (i Anm.\ til § 389 i
Encyklop.), siger han selv Intet til dets Opklaring, og hvad han herved
siger med Hensyn til de Philosopher, han her nævner, maa, især hvad
Cartesius og Leibnitz angaaer, forundre hver Den, som kjender disse
Tænkeres Philosophie. Men det er af flere Steder hos ham klart, at han ei
noksom kjender dertil.}, om Sjælens Udødelighed, om Guds Personlighed
o.\ fl., næsten urørte, ja ei engang har ladet Spørgsmaalet om
Philosophiens reale Mulighed, dens Natur og Væsen, samt om dens Forhold til
Empirien paa den ene, til Troen paa den anden Side, komme til nogen sand
indtrængende og gjennemtrængende Discussion. Naar en vis Tids Philosophie
begynder at sætte sig, for at udorganisere sig efter en vis bestemt
Grundtypus, og derover imidlertid en heel Menneskealder gaaer hen, vil dens
Bevægelsesperiode være forbi. Impulserne have da allerede virket, og
imidlertid har ogsaa Observationsaanden grebet om sig i Livets og i de
observerende Videnskabers Gebeet (men observerende ere de i Grunde alle),
saa at Forraad samle sig til Bedste for en ny Metamorphose i Philosophiens
Rige. Hos Hegel selv er et stort Forraad af slige Observationer at finde,
skjønt vistnok gjennemblandede med mange Konstlerier, og hans nærmeste
Virken for Fremtiden turde beroe paa, hvad Fremtidens philosophiske
Observatorer ville finde at tilegne sig og vide at føre sig til Nytte af
hans rige Forraad.

Hvad nu især den Hegelske Philosophies Forhold til Christendommen angaaer,
--- og derpaa maae vi da her især have Opmærksomheden henvendt --- saa er
vel nu allerede det Tidspunct indtraadt, da den Erkjendelse kan siges
noksom at have gjort sig gjeldende hos Mange, som ellers have følt sig
hentrukne til Hegel, at den Hegelske Christendomsphilosophie ikke svarer til
den ved Hegel selv til al Lære gjorte Hovedfordring, at den skal være
immanent, eller at den skal være Sagens egen sig udviklende Gang; hvoraf da
dog vel maa blive en Følge, at udaf Sagens eget Indre Videnskabens bevægende
Grundidee og det hele deri sig rørende Begreb maa gaae frem. Kun ved en sand
levende Vexelvirkning imellem den philosophiske Speculations Aand paa den
ene Side, og det givne Indbegreb, den givne Fylde af christelig Tro og
Erkjendelse paa den anden Side, kan Grundlaget for, og det sig heelt
igjennem Bevægende i, en Christendomsphilosophie hæve og klare sig frem. Kun
da kommer Christendomsphilosophien frem som et nyt Væsen, der har baade
Fader og Moder, og gjenudtrykker begges Naturer i en inderlig og væsentlig
Eenhed. Men den Hegelske Christendomsphilosophie, som vistnok er aprioriskt,
er ingenlunde immanent paa denne Maade. Den har ingenlunde udviklet sig udaf
en grundig philosophisk Fordybelse i selve Christendommen. Dens Forhold til
Læren om Udødeligheden, dens Trinitetslæres egentlige Beskaffenhed, hvorom
længere hen, saavelsom dens Miskjendelse af Følelsens og Troens egentlige
Væsen, samt det Suprematie, den, ligesom Fichte den Ældre, vil tillægge den
philosophiske Erkjendelse over Religionen, kan noksom lære Dette. Ikke at
tale om, at hos Hegel det hele Moralske ei bygges paa Religion, et Punct,
hvori vi atter see hele Tiden kommen ud over det Hegelske Stadium, skjønt
den hele Hegelske Philosophies besynderlige inverse Gang, som jeg pleier
kalde den, her har gjort Sit til at føre til denne store Ulempe. Som en
allerede færdig Philosophie træder den Hegelske Philosophie hen til
Christendommen, for nu og at tilegne sig den, der da maa lade sig arrangere
efter Principer, ved hvis Udvikling og Etablering den ikke retteligen og
fuldeligen har været paa Raad med. Christendommen har ei fra først af været
tilstede ved, og virksomt grebet ind i den Debat, hvoraf den hegelske
Philosophie har hævet sig frem. Spinozistiske Ideer have ganske anderledes
--- Enhver, som kjender til den ældre Schellingske Philosophie, og seer sig
om hos Hegel, maa kunne see det --- været i Bevægelse, da den Hegelske
Philosophie constituerede sig sit Grundlag. Med sine færdige
Grundanskuelser træder da Hegelianismen hen til det Christelige. Paa saadan
Maade kommer da Christendommen kun til at tjene den Hegelske Philosophie til
Føde; en paa Hegels Viis ontologiseret eller logisticeret Christendomslære
bliver da Resultatet, ingen sand speculativ christelig Theologie. At mange
enkelte Bemærkninger over det Christelige af immanent Udspring hos Hegel
kunne forekomme, skal ikke nægtes; jeg har allerede sagt, at der hos Hegel
er ikke Lidet, som peger hen ud over ham; men et forunderligt exoteriskt
Udseende faae ofte slige Yttringer hos ham. Men det maa herved ei heller
glemmes, at Hegels Philosophie hos ham selv allerede ganske havde
consolideret sig, eller fastnet sig, da de nye Bevægelser, det nye Røre fra
Grunden af, i de christne Erkjendelsers Verden brød frem med al Styrke.

Dog jeg maa her afbryde mig selv, for i Tide at lade en Hovedanmærkning faae
Plads, som jeg seer at være høist fornøden. Jeg seer nemlig, at man vil
kunne tage det Fremsatte, som om der til Grund derfor laae den Antagelse, at
Philosophien skal i de Regioner, hvori den kommer i Berøring med det givne
Christelige, gjøre dette, som saadant, til sit Grundlag. Dette er saa langt
fra at være min Mening, at jeg tvertimod i det Følgende, hvor jeg under en
særskilt Rubrik vil afhandle Philosophiens Begreb, vil søge at vise, at
Philosophien ingenlunde kan i Henseende til det Christelige bygges paa Tro,
som saadan, da Christendomsphilosophie kun kan komme frem ved den
speculative Idees Indtrængen ind i, og Gjennemtrængen hen igjennem, det hele
Christelige. Men et Andet er det, naar en Philosophie viser sig i høi Grad i
Strid med, eller i et underligt Misforhold til, de væsentligste christelige
Erkjendelser, og det endog saadanne, der høre til den almeen menneskelige
religiøse Idee- og Følelseskreds, hvilket gjelder om de ovenfor
(Pag.\ 294--97) omhandlede. Dette Misforhold maa da vidne imod hiin
Philosophie, men maa især blive af stor Virkning i saa Henseende i en Tid,
da hiin religiøse Idee- og Følelseskreds paa Ny er kommen i Bevægelse fra
Grunden af. Dette Misforhold kan da kun tages som et talende Beviis for, at
hiin den speculative Idees Indtrængen i det virkelige Christelige ei er
indtraadt, eller at den i denne Henseende ei har løst sin Opgave.

Forholdet er her, som i den Hegelske Naturphilosophie, hvis ikke faa og ikke
ubetydelige Misforhold til Alt, hvad Physik og Physiologie frembyder, vel og
kun behøve at fremhæves, for at vise, at denne Philosophie ei kan være egnet
til at anprises for Physikere og Physiologer, uden at man derfor vil
paastaae, at Naturphilosophie skal simpelhen bygges paa empirisk Physik og
Physiologie. Jeg vil saa meget heller her tage nogle Exempler heraf, som jeg
derved tillige vel kan befrie mig for den Nødvendighed i nærværende
Anmeldelse at røre videre ved dette Partie i Hegels Philosophie. Saaledes
som § 317 i Hegels Encyklopædie (baade i anden og tredie Udgave) lyder, kan
man, især naar man kjender til Hegels hele Naturphilosophie, og veed,
hvorledes man ofte maa søge hans Argumenter i enkelte Ord i hans Sætninger,
samt sammenligner § 317 med § 319, Anm.\ og § 320, ei Andet, end antage, at
Hegel anseer Gjennemsigtigheden grundet i en reen indre Homogeneitet i det
Materielle, samt at han troer, at enhver reen og fuldkommen udchrystalliseret
Chrystal maa være gjennemsigtig. Nu veed Enhver, som kjender Noget til slige
Ting, at Chrystallindsen i det menneskelige Øie, som vel Mere, end noget
Andet i Verden, er af den Beskaffenhed, at Lyset gaaer reent igjennem den,
og Mere, end noget Andet i Verden, maa kunne kaldes „ein Medium für das
Licht“, bestaaer af en Utallighed af Celler, opfyldte med Vædske, og at der
endda desuden i Øiet ere samlede flere Vædsker, gjennem hvilke Lyset her har
at gaae; saa at det er saare langt fra, at vi her have nogen Homogeneitet.
Fremdeles veed Enhver, som kjender Noget til Tingene, at Legemer, hvis hele
Natur er Gjennemsigtigheden imod, lige saa fuldt og lige saa fuldkomment
chrystallisere sig, som saadanne, hvis fuldeste Udchrystallisation ogsaa
fører til den reneste Gjennemsigtighed*)\footnote{Den anførte § 317 lyder i
den tredie Udgave saa: „In der gestalteten Körperlichkeit ist die
erste Bestimmung ihre mit sich identische Selbstischkeit, die abstracte
Selbstmanifestation ihrer als unbestimmter, einfacher Individualität ---
das Licht. Aber die Gestalt leuchtet als solche nicht, sondern diese
Eigenschaft ist (vorh.\ §) ein Verhältniß zum Lichte; 1) der Körper ist als
reiner Krystall in der vollkommenen Homogeneität seiner neutral-existirenden
inneren Individualisirung, durchsichtig und ein Medium für das Licht.“
Hertil føies som Anmærkning: „Was in Beziehung auf Durchsichtigkeit
die innere Cohäsionslosigkeit der Luft ist, ist im concreten Körper die
Homogeneität der in sich cohärenten und krystallisirten Gestalt. --- Der
individuelle Körper unbestimmt genommen ist freilich sowohl durchsichtig als
undurchsichtig, durchscheinend u.\ s.\ f. Aber die Durchsichtigkeit ist die
nächste, erste Bestimmung desselben als Krystalls, dessen physische
Homogeneität noch nicht weiter in sich besondert und vertieft ist.“}.

Jeg tager et andet Exempel. I § 354 gjennemgaaes Sensibilitet, Irritabilitet
og Reproduction saaledes, at ved hver af de to første angives tre Momenter,
svarende til denne Tredeling selv. Som det Moment i Sensibilitetens System,
der nu især skal udgjøre Sensibilitetens Moment deri, nævnes da
Knokkelsystemet, hvorimod Hjernen med de Nerver, som høre til Cerebral- og
Dorsalsystemet, og det, som det udtrykkeligen siges, lige saa vel Nerverne
for Fornemmelsen, som for Bevægelsen, skulle høre til Det, der i
Sensibilitetens System skal henregnes til Irritabilitetens Moment.
Forundringen herover kan ikke formindskes, naar man seer hen til det,
forresten høist underlige, snevre Begreb, Hegel giver om Sensibilitet i
Almindelighed i § 353; Ordene lyde i tredie Udgave (med nogen Forskjellighed
fra deres Lydende i den af mig ellers meest benyttede anden Udgave)
saaledes: „Das animalische Subject ist daher α) sein einfaches, allgemeines
Insichseyn in seiner Aeußerlichkeit, wodurch die wirkliche Bestimmtheit
unmittelbar als Besonderheit in das Allgemeine aufgenommen und dieses in ihr
ungetrennte Identität des Subjects mit sich selbst ist --- Sensibilität.“

Til Irritabiliteten finde vi siden hele Blodløbet henregnet, hvorom det da
tilsidst hedder, at det er „ein allgemeines Preisgeben an die Reproduction
der übrigen*)\footnote{Dette „übrigen“ har ifølge Sammenhængen Hensyn til
Hjertet, som da maatte synes ikke at skulle selv faae nogen Andeel i hiin
Reproduction. Men ved en saadan liden \textit{incuria} i Udtrykket, vilde
det være høist ubilligt at hænge sig. Jeg har her kun hensat denne
Anmærkning, fordi den opmærksomme Læser naturligviis vil spørge, hvortil
dette „übrigen“ sigter.} Glieder, daß sie aus dem Blute sich ihre Nahrung
nehmen.“ Deraf bliver da en Følge, at, naar vi nu komme til det tredie
Hovedled, selve Reproductionen, saa kommer den egentlige Reproduction ikke
til at indtage denne tredie Hovedplads, og det uagtet Reproductionen i § 353
opføres som høieste Led i Treheden, og kaldes „die Einheit dieser Momente“
(nemlig af Sensibilitet og Irritabilitet), „die
negative**)\footnote{Jeg bør her ikke lade dette Tillægsord: negative, gaae
uforklaret hen. Hegel lader overalt i sine Treheder et relativt Første træde
ud i sit Andet, som i sin Modsætning tillige, men derigjennem, ved denne sin
Modsætnings (eller sit Negatives) Negation gaae tilbage i sig selv igjen, og
derved træde frem i en større Fylde, som det mere Concrete. Denne
Negationens Negation kalder han i Særdeleshed Negativiteten. Men her paa
dette Sted er den rigtignok kun at spore i Ordet.} Rückkehr zu sich selbst
aus dem Verhältnisse der Aeußerlichkeit und dadurch Erzeugung und Setzen
seiner als eines Einzelnen.“ Den hele tredie Hovedplads, hvor da Det, som nu
her er angivet, skulde søges, kommer da, eftersom Dette allerede forud er
borttaget og bragt an under Irritabiliteten, kun til at indtages af
Fordøielsessystemet og af Lymphasecretionen, altsaa af de Systemer, hvorved
Det beredes, der skal føres over i Blodet, for derfra at føres til de
Puncter, hvor den egentlige Reproduction først gaaer for sig, saa at ikke
denne selv, men kun dens Præliminarier, saa at sige, opføres for os paa den
Plads, hvor vi skulde finde hiin „Rückkehr zu sich selbst“ og hiin
„Erzeugung seiner als eines Einzelnen.“

Naar man nu, ved at anføre disse og mange flere aabenbare Misforhold, hvori
Hegels Naturphilosophie staaer til Physik og Physiologie, med det meget
deels \emph{Phantastiske}, deels \emph{Nebuløse}, deels \emph{Tautologiske},
deels om en høist begrændset Naturanskuelse
Vidnende*)\footnote{Til Exempel paa alle disse fire Ting kan hans § 275
tjene: „Die erste qualificirte Materie ist sie (nemlig Materien), als ihre
reine Identität mit sich als Einheit der Reflexion-in-sich, somit die erste,
selbst noch abstracte Manifestation. In der Natur daseyend ist sie die
Beziehung auf sich als selbstständig gegen die andere Bestimmungen der
Totalität. Dieß existirende allgemeine Selbst der Materie ist das Licht, ---
als Individualität, der Stern, und derselbe als Moment einer Totalität, die
Sonne.“ Denne Paragraph staaer som første § under Rubriken: „die freien
physischen Körper.“}, man i Hegels Naturphilosophie møder, som Beviser for
det høist Utidige i at ville anprise Hegels Philosophie for philosopherende
Physikere og Physiologer, ja at ville tage saadanne det ilde op, at de ville
ud over en slig Naturphilosophie, der, --- om den nogensinde kan have havt
et Slags Betydning, som et Slags Forsøg paa at philosophere i
Naturvidenskaben --- dog nu maa kaldes aldeles forældet, vil man da derved
afgjøre Spørgsmaalet om Forholdet imellem Philosophie og Empirie derhen, at
hiin med Hensyn til Erfaringen skulde bygge paa Empirie? Man vil kun fordre
anerkjendt, at den skal kunne staae sin Prøve med Hensyn til Erfaringen, og
derved skal vise, om det er den sande, virkelige Natur, hvis Philosophie man
har faaet. Og saaledes maae vi, med Hensyn til Philosophiens Forhold til
Christendommen, paastaae, at den vel ingenlunde skal, som Philosophie, bygge
paa Christendommens Indhold, som givet, men at den dog meget slet staaer sin
Prøve med Hensyn til denne, naar den ei har kunnet komme bedre tilrette med
saa vigtige Puncter, som de, jeg forhen har nævnet, end den har gjort det.
Nu veed jeg jo vel, at der er den Forskjel imellem Naturbetragtning og
Christendom, at, medens i hiin Ingen kan afvise, hvad Physikeren empiriskt
kan paavise eller godtgjøre, saa kan man i den religiøse Ideekreds erklære
alt Det, man ei kan komme tilrette med, eller ei vil vide af, for ureelt,
eller for underordnet. Men derom skal jeg tale nærmere, naar jeg kommer til
at afhandle Philosophiens Begreb i det Hele. Hvad jeg her har for, er kun at
betragte Hegel i Forhold til vor Tid.

Træder man nu i vor Tid netop især hen til Saadanne, som med ungdommelig
Varme og Kraft have været ret levende inde med i hine Christendommens nye
Foraarsbevægelser, og det maa man dog antages at gjøre, naar man nu i vor
Tid vil stille sig lige overfor yngre Theologer, som man tænker sig med
speculativ Aand, saa falder dog Opfordringen til at kaste sig ganske ind i
den Hegelske Philosophies Studium, som en underlig Revocation, eller som en
Anmodning om, efterat Skridt ere skeete hen til en meget avanceret
Operationslinie, at drage tilbage til en vis ældre Operationslinie, fordi
denne holdes for at være bedre befæstet, uagtet man ved hiin har aabnet sig
en mere indtrængende Adgang til Hjertet af det forjættede Land, som paa Ny
skal indtages. Naar hos yngre Theologer, af bestemt og dygtig philosophisk
Drivt og Aandsgave, Spørgsmaalene om Forholdet imellem Tro og Philosophie
reise sig paa Ny, som de, der ville arbeide hen til et andet Grundlag for en
speculativ Dogmatik, end et saadant, som Hegel kan frembyde --- hvad er det
saa i Grunden Andet, end at de bemærke, at man, for at faae en ret
tilfredsstillende Christendomsphilosophie, ikke en aprioriske, men en
virkelig immanent, først maa være fuldkomment indlevet i Christendemmens
indre Gehalt, og have ladet den selv raade og selv klare sig til
Erkjendelse, medens de vel nok maae see, at den Hegelske Philosophie ikke
har en saadan sand Intussusception af Christendommen i sit Grundlag? Man kan
vel vide, at det maa være mig, da jeg i mange Aar har, baade i Skrivter og
ved mundtlige Foredrag, især ved de sidste, arbeidet for en
Christendomsphilosophie af saadan Art, meget kjært, at nu Følelsen for, at
den kun ad saadanne Veie er at finde, gjør sig gjeldende selv hos Saadanne,
som en Tid lang have levet sig ind i Hegel, der iøvrigt vistnok ikke kan
Andet, end have noget meget Tiltrækkende for speculative Hoveder, og der
tillige vel nok en Tid lang, indtil man faaer Impuls til, engang at see sig
videre om i Verden og forfare, om der boer Nogen bag de blaae Bjerge i
Horizonten, kan synes En at have arrangeret Alt ypperligen, eller i al Fald
leveret en bevægelig og betydningsfuld Grundtypus at arrangere sig efter.
Det er da at haabe, at det vil vise sig, at det i Christendommens Rige
kommer an paa noget ganske Andet, end hvilkensomhelst Ontologisering, som
og, at den christelige Fyldes Rige er noget Andet og Mere, end ein Reich des
Wissens, eller end et Rige, som skulde finde sit βασιλεία τῶν οὐρανῶν i en
Philosophie.

„Naturwissenschaft a priori ist der Tod aller Naturphilosophie“ har Steffens
engang sagt. Jeg kan ei mere angive, hvor hos Steffens denne Yttring staaer;
jeg skulde troe, det er i hans „über die falsche Theologie u.\ s.\ w.,“ men
jeg har ikke Bogen ved Haanden, for at kunne see efter. Des vissere veed jeg,
at jeg hos ham har fundet disse --- ved den Stilling, som Ordet
Naturphilosophie her har faaet, mærkelige --- Ord; thi jeg har gjemt dem hos
mig fra den Tid, jeg læste dem, og er ofte vendt tilbage til dem. Jeg maatte
vel mindes dem nu; thi i Sandhed en Dogmatik i den Hegelske Form og efter
den Hegelske Aand, tør vel foranledige til at gjøre den Erindring, at, naar
Christendomsphilosophien ikke udaf en speculativ Indtrængen i
Christendommens eget Dyb udvikler sig, og det ei blot i Henseende til sin
Gehalt, men og i Henseende til sin Organisation, men skal lade sig modulere
efter en forud medbragt Ontologies Typus, saa mortificerer man al sand
Christendomsphilosophie; og naar man vil kalde Sligt speculativ
Christendomslære, speculativ Dogmatik, saa veed jeg ikke, hvor man er henne
med sine Benævnelser. Christendommen forekommer mig da dog at være i langt
bedre Hænder i et Forsøg, som Dr.\ Rothes, end i en Dogmatisering af hiin
Art; thi, naar en saadan gaaer ud paa at hegelianisere Christendommen, saa
gaaer Rothes Forsøg ei ud paa noget Lignende, men er fra først af i Grunden
gaaet ud paa, at christianisere Naturen, og attraaede snarere at arbeide hen
paa at levere Bidrag til en christelig Naturbetragtning, end til en
speculativ Christendomsphilosophie. Rothes Forsøg er i Grunden \emph{hans}
Treenighedens gjentagne Aabenbaring i den hele store Natur. Jeg troer
virkelig ved at sige Dette, at have sagt Noget, som turde betegne det rette
Forhold i Rothes Speculation i det Hele taget. Man seer, at han allerførst
spørger sig: hvad lærer Christendommen Dig? Dernæst spørger han sig, om han
nu gjenfinder det Samme overalt, hvor han seer sig om i Naturens og den
indre Mennesketilværelses Rige. Dermed hænger det ogsaa sammen, at han ei
gjør Fordring paa det Navn Philosophie for sine Betragtninger, idet ogsaa
han synes at have den Tanke, at Philosophie skal være reen apriorisk,
hvorimod han erklærer sine Betragtninger at beroe paa Væsensskue, og denne
kalder han nu Speculation. Jeg kommer naturligviis til paa sit Sted at
omhandle Prof.\ Heibergs Angreb paa Rothe i denne Henseende. Men skjønt
Rothes Christendomsphilosophie langt snarere er at kalde immanent, d.\ e.\
hjemmevoxet paa Christendommens egen Mark, end nogen Hegelsk er det, saa er
dog hans Philosophering ei naaet til at blive virkelig immanent; thi, hvad
han i Naturens Rige er bleven vaer, som svarende til det Christelige, saavidt
hans Syn naaede, har havt en altfor stor Tilbagevirkning paa Udtydningen af
de christelige Lærdomme selv, med alt det at hans Agt var at tage disse og
holde paa disse, som den hellige Auctoritet gav ham dem; thi sin præstelige
Charakteer fornægter har intet Øieblik, og det er høist paafaldende, at
Prof.\ Heiberg ofte synes reent at glemme, at han har en christelig Præst for
sig, for hvem saadanne Steder, som f.\ Ex.\ Pauli Ord i 1 Cor.\ 13, ---
hvor der tales om den vor Erkjendelse, som her blot er Stykværk, og om den
Skuen Ansigt til Ansigt, hvilken vi forvente --- staae som lysende Stjerner.

Paa Veien til en sand Christendomsphilosophie kommer man da vist ikke ved at
lade sig hegelianisere, og derpaa, som en vel hegelianiseret Philosoph, at
ville see sig om i Christendommens Enemærker. Den rette Vei til at komme til
Christendomsphilosophie kan kun være, at man ogsaa i philosophisk Henseende
lader sig christne, og ikke med et fremmed Øie, men med det christelige
Sindelags eget Øie seer, hvad man i Christendommen har, idet man kun vil
tale om, hvad man har oplevet; omtrent som Veien til at blive Naturphilosoph
maa blive den, at man først lader sig gjøre til Physiker og Physiolog. At en
sand Naturphilosophie kun kan komme i Stand ved et grundigt philosophiskt
(ikke blot empiriskt) Studium af selve Naturen, er vel uimodsigeligt. Hegel,
hvis Naturphilosophie udgjør hans Philosophies \textit{partie honteuse},
skjønt ingenlunde dets \textit{partie modeste}, kan være et stort Exempel
til Advarsel i denne Henseende, eller til Advarsel mod en ikke immanent, men
efter en fremmed Ontologies Tilsagn og Indskydelser schematiseret
Naturphilosophie. Som der nu hører Fortrolighed med Naturen og en solid
Skole, man hos den selv har gjennemgaaet, til at naae saavidt, at man
betragter Naturen retteligen med et philosophiskt Øie, og bliver vaer, hvad
den selv frembyder for det speculative Blik, til Tilfredsstillelse for den
speculative Trang og Begjær, saaledes maa man og respirere i Christendommens
og den christelige Troes eiendommelige Luftninger, og røre sig i dens egen
Fylde, for at komme til en sand speculativ Opklaring i Christendommens
Regioner. Det er Det, Anselmus viser hen paa i de noksom bekjendte Ord,
hvilke Schleiermacher har sat som Motto paa sin Dogmatiks første Deel, hvori
det hedder: \textit{qvi non crediderit, non experietur, et qvi expertus non
fuerit, non intelliget.} Troen staaer i samme Forhold til Philosophien,
hvori Sagen selv staaer til den philosophiske Betragtning af den; thi Troen
er Christendommen selv i virkelig Tilstedeværelse og Leven i Gemytet.
Speculativ Dogmatik derimod er ikke selve Christendommen. Jeg tillader mig
herved at minde om, hvad jeg flere for Aar siden engang har
sagt*)\footnote{See mit Philosoph.\ Archiv første Hefte Pag.\ 62.}: „at, om
end Philosophien er Aanden i den christelige Erkjenden, saa er dog den
tillidsfulde, ud over al Seen og Erkjenden gaaende, alt det Erkjendte og
meget Mere, end det, som i en Livskjærne i sig indsluttende Tro, Sjælen i den
christelige Existents.“

I Sandhed, jo mere jeg er gaaet ind i den Hegelske Philosophie, saaledes som
den ligger for os i de af Hegel selv udgivne Hovedværker: hans Logik i tre
Dele, hans Encyklopædie, hvori vi have faaet hans hele Systems Grundtræk,
flere Gange gjennemarbeidede, i tre Udgaver, og endeligen den specielle
Udførelse af hans Lære om den af ham saakaldte objective Aand, hvilken vi
have faaet under Navn af Philosophie des Rechts, --- desto mere har jeg
vistnok maattet erkjende og med Fornøielse sande, at den bærer en speculativ
Gehalt i sig, som ei kan Andet, end udøve en mægtig Tiltrækning paa
speculative Hoveder, naar de først have faaet Brud paa den, hvortil hos os
uden Tvivl Professor Heiberg har forhjulpen Adskillige. Jeg skulde maaskee
ikke let kunne angive noget Værk, hvis Studium, ved Siden af de platoniske
Dialogers Studium, mere kunde være dannende og udviklende for speculativ
Tænkning i logisk Henseende, end Hegels Logik, samt eet og andet Partie i
hans Aandsphilosophie, især Alt, hvori han forfølger den autonomiske
Proces*)\footnote{Herved mener jeg hans Skildring af Gangen i hen hele
Maade, hvorpaa den individuelle Bevidsthed tager sig selv i Besiddelse, idet
den tager sit Livs Indhold i Besiddelse, og ligesom overmander det, men med
det Samme ogsaa tages i Besiddelse af den objective Aand. Det herhen hørende
Psychologiske er ei allene det Fremherskende, ja næsten det ene Herskende i
hans Psychologie, men udgjør ogsaa noget særdeles Tiltrækkende i den første
Afdeling af hans Philosophie des Rechts, eller den Afdeling deri, som udgjør
den egentlige Naturret, der forresten er ubetydelig, som man da snart
mærker, at den ei er af en Forfatter, som er kommen ind i de herhen hørende
Debatter. (Langt bedre, og i Sandhed fortrinlig, er hans om en temmelig
frisindig Tænkemaade vidnende Statslære, skjønt man ogsaa her snart bliver
vaer, at Forfatteren ei er Jurist.)}, om jeg saa maa kalde den. (Hans
ligesaa barokke, som \textit{affreuse} Naturphilosophie maa blive udenfor).
Men jo mere jeg har skuet ind i hans hele Philosophie, desto mere har det
ogsaa maattet være mig underligt, at Mænd med frit Omløb i Tankeriget ikke
have kunnet blive vaer, hvor utilfredsstillende dog i mangen Henseende denne
Philosophie ikke kan Andet end være for Den, der kjender til Livet og dets
Indhold, og kjender til hvad der ellers i Philosophien er traadt frem, endog
i selve den Schellingske Kreds, i hvilken sidste Henseende jeg kun vil nævne
Steffens og Franz Baader, samt er skredet nogenlunde frem med Tiden, især i
de christelige Overveielsers Sphære.

Selve den Hegelske Philosophie indeholder Nok til at være en Stagnation i
denne Henseende imod; den indeholder Nok, der peger ud over dens eget
Stadium, Nok for at afholde fra den Tanke, at her kunde man nu indtil videre
slaae sig til Ro, og give sig hen i Hvilens Fylde. At man ret snart bliver
Dette almindeligen vaer, er ogsaa for de Hegelske Bestræbelsers egen Skyld
saare ønskeligt. Man har, hvad Virkningen af den Hegelske Philosophies
Studium i Tydskland angaaer, villet bemærke, at det synes at Hegel blandt
sine Modstandere besidder nogle af sine bedste Læsere, hos hvilke hans
Philosophie meest har baaret Frugt. Dette kan ei forundre, og det er for
dens egen Skyld kun at ønske, at det underlige unaturlige Forhold maa
ophøre, ifølge hvilket der har dannet sig en Tilstand, at det seer ud, som
om det kom an paa, enten at bekæmpe den, som en Philosophie, man aldeles
maatte lade falde, eller at forfægte den, som en Philosophie, der i Eet og
Alt var saa tilfredsstillende, at det kun kom an paa at udbrede den og gjøre
den gjeldende overalt. Saa synes at blive vaer, at, for retteligen at nytte
den her meddeelte store Gehalt, ja jeg kan sige, for tilbørligen at lade den
vederfares Ret, ved at lade den vorde saa frugtbringende, som den kan blive
det, maa man netop gaae ud over den. Istedetfor at optage den som et stort,
høist priisværdigt Bidrag til den philosophiske Stræbens og det
philosophiske Værks Fremme, dele sig de Philosopherende for en stor Deel,
forsaavidt de paa en livfuld Maade indlade sig med Hegel, i Modstandere, som
blot forholde sig angribende, Tilhængere, som blot forfægte og yttre sig
paaholdende, og Saadanne, som ønske at nyde Godt af hvad her er givet, uden
at ville slaae sig til noget Partie, men som dog ei kunne undgaae at drages
ind med i denne Antagonismus. Til de sidste turde nærværende Recensent vel
være at henregne. Saa at, tvertimod den hele Hegelske Philosophies egen
Synsmaade i Henseende til Det, som i Philosophiens Historie er traadt stærkt
indgribende frem, --- hvorefter man ei skal gaae ud fra, at et philosophiskt
System netop enten skal antages for gyldigt, eller gjendrives som ugyldigt,
--- trækkes nu Hegel selv med ind under et saadant \textit{aut-aut}, maaskee
til Straf, fordi han selv med Hensyn til Andre, især Jacobi, saa lidet har
havt for Øie, hvad han (Logik III, strax i Beg., Pag.\ 8 i den første Udg.)
fordrer med særdeles Hensyn til Spinoza.*)\footnote{„Die wahrhafte
Widerlegung“ --- saa hedder det her --- „muß in die Kraft des Gegners
eingehen und sich in den Umkreis seiner Stärke stellen; ihn ausserhalb
seiner selbst angreifen, und da Recht zu behalten, wo er nicht ist, fördert
die Sache nicht.“ --- „Die einzige Widerlegung“ --- lægges der til --- „des
Spinozismus kann daher nur darin bestehen, daß sein Standpunct zuerst als
wesentlich und nothwendig anerkannt werde, daß aber zweitens dieser
Standpunct aus sich selbst auf den höhern gehoben werde.“ --- At saadan Art
af Modvirkning allene skulde være den rette, kan iøvrigt vist Ingen
indrømme; thi megen Modvirkning maa dog have sit Grundlag i, at man friere
og i en større Omkreds seer sig om i Livet. Men Hegel selv kunde indbyde
til, at Nogen netop behandlede ham, som han her har fordret Spinoza
behandlet, og selv har handlet med Hensyn til Spinoza (hvorved denne
forresten vistnok alt for meget er bleven ham hans egentlige Grundlag).}
Dertil have da de altfor eensidige, blot og bart Hegelske Tilhængere, samt
tillige Hegels egne store Eensidigheder, saavelsom hans slemme,
sædvanligviis haanende, absolutistiske Maade at polemisere paa, bidraget
meget, især da denne Polemik har smittet mange af Tilhængerne, deels med den
Fornemhed, deels med den Sufficence, som Eensidighed dobbelt kommer til,
naar den, under en stor Auctoritets Beskyttelse, troer sig berettiget til
haanlig at afvise. Hvoraf Følgen atter er bleven, at Modstanderne dobbelt
ere blevne æggede, og have viist, at ei blot Kjærlighed, men ogsaa
Indignation gjør skarpseende, kun paa forskjellig Maade. --- Men ligesom
neppe Den, der allene betragter Hegels Philosophie med et mod samme indtaget
Sinds Øie, vil kunne blive fuldkomment vaer, hvorledes man rettest fører ud
over den, og derved bedst overvinder den, saaledes vil ogsaa neppe Den med
nogen sand Frihed, altsaa paa den rette frugtbringende Maade, tilegne sig
det Meget, som her er givet, der ikke tillige kan see, at Meget af det
Betydeligste i denne Philosophie staaer i et aabenbart Misforhold, deels til
Philosophiens egentlige Væsen, deels til Livets Indhold, ogsaa%
% [text to be added: pp. 33-- (PDF 44--), continue in ~10-page batches]

\end{document}
